\section{Auditoría de fuentes y pregunta central del curso}

La pregunta central de esta clase no es «si importa el número de parámetros», sino: cuando un modelo de lenguaje ya genera respuestas con fluidez, ¿qué mecanismos adicionales hacen falta para convertirlo en un razonador, un modelo de conocimiento o un componente de juicio ético más confiable? Yejin Choi responde con tres trabajos de investigación: Maieutic Prompting, Symbolic Knowledge Distillation y Delphi. Los tres dependen de modelos de lenguaje grandes, pero todos añaden fuera del modelo estructura, datos, un critic o restricciones. Por eso no pueden reducirse a la frase «un modelo pequeño vence a uno grande».

\subsection{Cómo se reconstruyó esta clase}

El sitio oficial del curso publica la grabación y los tres artículos asignados, pero no un slide deck independiente. Esta clase toma como fuente visual principal la grabación oficial de Stanford Online en 1080p: un muestreo cada dos segundos produjo 2,253 fotogramas, la herramienta de fotogramas estables generó 207 candidatos de alta exhaustividad y una depuración manual de duplicados dejó 56 estados didácticos independientes. Todas las imágenes finales se extrajeron del cuadro completo de 1920x1080. Después del minuto 52:40, la sesión de preguntas vuelve muchas veces a diapositivas anteriores, así que las imágenes repetidas no se insertan de nuevo; las explicaciones orales se conservan mediante los subtítulos oficiales \texttt{en-US} hechos por personas.

\begin{longtable}{@{}L{0.18\textwidth}L{0.27\textwidth}L{0.43\textwidth}@{}}
\toprule
Nivel de evidencia & Material local & Función en los apuntes \\
\midrule
Grabación oficial & Video de 1:15:05 y 56 teaching states & Determina el orden de la clase, las diapositivas de fórmulas, las tablas de resultados, las demostraciones públicas, las respuestas a la controversia y la arquitectura hybrid \\
Subtítulos oficiales hechos por personas & \texttt{lecture16.en.srt} y timed transcript & Recuperan las motivaciones de la expositora, los caveats de las comparaciones, los sesgos de los datos, el costo de los sistemas y la Q\&A \\
Artículos asignados & Maieutic, Symbolic KD, Delphi & Verifican algoritmos, fórmulas, alcance experimental, composición de los datos y limitaciones \\
Artículos de contexto & COMET, COMET-ATOMIC 2020 & Explican la relación entre el symbolic graph y el neural knowledge model \\
Archivos de auditoría & manifest, coverage, ledger, blueprint & Demuestran que cada nodo visual y de teacher voice tiene un destino didáctico explícito \\
\bottomrule
\end{longtable}

\begin{warningbox}{Límite histórico: 14 de febrero de 2023}
Esta clase trata los sistemas y debates de cuando ChatGPT llevaba apenas unos meses disponible al público. Los reasoning models, los métodos de post-training y los productos de machine ethics posteriores no pueden proyectarse hacia atrás como evidencia de la clase. En particular, Delphi hybrid seguía siendo un work in progress en el momento de la clase y no debe describirse como un sistema maduro ya desplegado.
\end{warningbox}

\subsection{David y Goliat: el título no es un manifiesto contra el Scaling}

La diapositiva del título compara los modelos enormes con Goliat y los sistemas más pequeños pero estructurados con David. Al leer esta figura hay que conservar los matices: la expositora reconoce que los modelos grandes son superiores en muchas tareas; lo que rechaza es concluir que «el sentido común ya está resuelto» a partir de unos pocos ejemplos fluidos. Lo que busca esta clase es otra vía de mejora: evaluation más realista, task mejor focalizadas, datos de mayor calidad y algoritmos capaces de verificar la salida de un modelo de lenguaje.

\lecturefigure{slide-01-title-david-vs-goliath.jpg}{David vs. Goliath: el arte de los rankings en la era de los modelos de lenguaje enormes}{Diapositiva V001 recuperada de la grabación oficial; Stanford CS25 Lecture 16}

La siguiente slide toma de «El arte de la guerra» (孙子兵法) una división del método de investigación en tres pasos: primero conocer la distribución de los fallos, después elegir el campo de batalla que realmente discrimina y por último diseñar armas nuevas. No es un adorno literario, sino la evidence discipline de toda la clase: si el benchmark solo recompensa un dataset artifact, ni la puntuación más alta demuestra que el modelo domina la underlying task.

\lecturefigure{slide-02-art-of-war-lessons.jpg}{Know your enemy, choose your battles, innovate your weapons}{Diapositiva V002 recuperada de la grabación oficial; Stanford CS25 Lecture 16}

\teachervoice{La expositora usa con ironía la expresión «existential crisis» para describir el impacto de ChatGPT, pero enseguida replantea el problema como una hasty generalization: responder bien una vez una pregunta de Winograd no equivale a tener un sentido común físico estable. Lo que repite una y otra vez no es que «la escala no sirve», sino «smaller can be better» y «knowledge is power».}

\subsection{Glosario: distinguir primero los términos afines}

Más adelante aparecerán con frecuencia language model, knowledge model, commonsense graph, critic, descriptive morality y normative ethics. Si no se distinguen antes sus objetos, es fácil confundir «generar una frase razonable», «poseer conocimiento reutilizable» y «emitir un juicio ético» como si fueran la misma capacidad. La tabla siguiente da las definiciones de trabajo que usa esta clase.

\begin{longtable}{@{}L{0.19\textwidth}L{0.27\textwidth}L{0.42\textwidth}@{}}
\toprule
Término & Problema que resuelve & Mecanismo central y relación con el curso \\
\midrule
Language model & Cómo se distribuye el siguiente token & Aprende $p(x_t\mid x_{<t})$ a partir de estadísticas del texto; la fluidez puede estar relacionada con el conocimiento del mundo, pero no equivale a él \\
Knowledge model & Qué se puede inferir de una entidad, un evento o una relación & Conserva conocimiento reutilizable mediante typed relations, interfaces de recuperación u objetivos de generación especializados \\
Commonsense & Conocimiento práctico y razonamiento ampliamente compartidos en situaciones cotidianas & Abarca lo físico, lo social, lo causal, las intenciones y las reacciones; no es idéntico en todas las culturas \\
Symbolic graph & Cómo representar de forma explícita nodos y relaciones & ATOMIC usa nodos de eventos y relation types para formar una estructura de conocimiento verificable \\
Critic & Qué contenido generado por la máquina debe conservarse & Aprende a discriminar la calidad y filtra las teacher samples; el propio critic también se equivoca \\
Descriptive ethics & Cómo suelen juzgar las personas & Delphi predice sampled human judgments; no pretende dar la respuesta correcta definitiva \\
Normative ethics & Cómo deberían juzgar las personas & Requiere principios, argumentación y participación institucional; no puede derivarse automáticamente de las etiquetas mayoritarias \\
Reflective equilibrium & Cómo conciliar las intuiciones sobre casos con los principios abstractos & Corrige de forma iterativa entre el bottom-up judgment y el top-down constraint \\
Adversarial test & Si el sistema se mantiene estable ante entradas anómalas, compuestas o construidas deliberadamente & Sirve para exponer la brecha entre dataset success y task understanding \\
\bottomrule
\end{longtable}

\subsection{Cuatro preguntas que atraviesan toda la clase}

Al leer los tres sistemas que siguen, conviene plantear de manera constante cuatro preguntas. Primera: ¿el base LM aporta fluent generation, latent knowledge o calibrated truth? Segunda: ¿qué estructura coloca el sistema fuera del modelo: logic, knowledge graph, critic o moral principle? Tercera: ¿los resultados miden realmente accuracy, consistency, human plausibility o deployment safety? Cuarta: ¿los fallos pueden localizarse y corregirse, o solo se obtiene otra puntuación global no interpretable?

\begin{importantbox}{La Recipe común de esta clase}
Primero, el modelo grande genera candidatos; después, una estructura explícita o un modelo especializado reduce el espacio de candidatos; por último, el resultado se usa para entrenar un sistema más pequeño y especializado. Maieutic restringe las explanations con logic, Symbolic KD restringe el knowledge corpus con un critic y Delphi / hybrid restringe el judgment con commonsense data y moral constraints.
\end{importantbox}

\subsection{Resumen de la sección}

La auditoría de fuentes muestra que esta clase trata sobre la división de funciones en un sistema y los límites de la evidencia. David-versus-Goliath no niega la escala; exige colocar el modelo de lenguaje dentro de un pipeline verificable. Cada sección posterior debe explicar de dónde provienen los candidatos, quién impone las restricciones, qué demuestran las métricas y qué no demuestran.
